\section{Introdução}
A recuperação de informação (RI) moderna raramente depende de um único sinal. Modelos lexicais como o BM25 são rápidos, interpretáveis e fortes quando consulta e documento compartilham termos, mas falham diante de sinonímia e variação de vocabulário; \emph{bi-encoders} densos comparam significados e toleram paráfrases, porém podem diluir termos técnicos raros e truncar textos longos; \emph{cross-encoders} avaliam consulta e documento em conjunto, com mais precisão e mais custo; e a geração aumentada por recuperação (\emph{Retrieval-Augmented Generation}, RAG) usa os documentos recuperados como evidência para uma LLM responder \cite{lewis2020rag}.

Este trabalho implementa e avalia experimentalmente o pipeline completo \emph{documentos $\rightarrow$ pré-processamento e indexação $\rightarrow$ recuperação lexical + semântica $\rightarrow$ busca híbrida $\rightarrow$ reranking $\rightarrow$ seleção de contexto $\rightarrow$ LLM $\rightarrow$ resposta} sobre o \textbf{NFCorpus} na versão do benchmark BEIR \cite{boteva2016nfcorpus,thakur2021beir}: \NDocs{} documentos médicos (título e resumo) e \NQueries{} consultas derivadas do NutritionFacts.org, com julgamentos de relevância graduados (0--2) e, em média, \AvgRelTest{} documentos relevantes por consulta no \emph{test}.

\textbf{Objetivos.} (i) construir do zero um índice invertido e a busca lexical BM25; (ii) tornar o sistema tolerante a erros ortográficos (k-gramas, Jaccard e distância de edição) e medir quanto do desempenho perdido é recuperado; (iii) recuperar semanticamente com um \emph{bi-encoder} pré-treinado; (iv) fundir os rankings com \emph{Reciprocal Rank Fusion} (RRF); (v) reordenar o Top-50 com um \emph{cross-encoder}; (vi) avaliar o efeito da recuperação na qualidade das respostas de uma LLM (RAG $\times$ sem RAG). As conclusões se apoiam em avaliação quantitativa, análise por consulta e investigação de casos de sucesso e falha, sempre com valores medidos nos próprios experimentos. Todas as escolhas de modelos e parâmetros foram feitas no conjunto \emph{dev}; o \emph{test} só foi usado na avaliação final.

\section{Metodologia}
\subsection{Pré-processamento}
Cada documento é representado pela concatenação \emph{título + espaço + texto} em \emph{todos} os sistemas (lexical, semântico e reranking). Para o índice, o texto é (1) convertido para minúsculas e sem acentos (NFKD), (2) tokenizado em sequências alfanuméricas (\texttt{[a-z0-9]+}; hífens e apóstrofos separam tokens), (3) filtrado pela lista de \emph{stopwords} do inglês do NLTK \cite{bird2009nltk} e (4) reduzido ao radical pelo stemmer \PreStem{} do NLTK. Consultas passam exatamente pelo mesmo pipeline. As etapas (stopwords e stemming) foram escolhidas no \emph{dev} (Tabela~\ref{tab:dev_pre}); a lematização não foi avaliada. Uma segunda visão do texto, os \emph{tokens normalizados antes de stopwords e stemming}, forma o dicionário da correção ortográfica (\DictSize{} termos). Bibliotecas externas: \texttt{unicodedata}/\texttt{re} (biblioteca padrão) e NLTK \NltkVer{}.

\subsection{Índice invertido (implementado do zero)}
O dicionário é um \texttt{dict} Python que mapeia cada termo a uma \emph{posting list} formada por dois vetores de inteiros (\texttt{array('I')}) paralelos e ordenados por documento: identificadores internos e TF. O DF do termo é o comprimento da lista. Guardamos ainda o tamanho $|d|$ de cada documento (tokens após o pré-processamento) e as estatísticas globais ($N$, comprimento médio, vocabulário, total de \emph{postings}). A consulta é processada \emph{termo a termo} (\emph{term-at-a-time}): para cada termo da consulta percorre-se apenas a sua posting list, somando as contribuições em um acumulador por documento; nunca se varre a coleção. O Top-$k$ sai de uma seleção por \texttt{heapq} sobre os acumuladores, com desempate pelo menor identificador. A Tabela~\ref{tab:index} traz as estatísticas do índice e os tempos de construção e de consulta.

\begin{table}[t]
\centering
\caption{Índice invertido construído do zero.}
\label{tab:index}
\small
\begin{tabular}{lr}
\toprule
Estatística & Valor \\
\midrule
Documentos & 3.633 \\
Termos do vocabulário & 18.958 \\
Total de postings & 359.439 \\
Tokens indexados & 602.027 \\
Comprimento médio do documento & 165,7 \\
Construção (pré-proc. + estruturas) & 0,82 s + 0,22 s \\
Tempo médio por consulta & 0,25 ms \\
\bottomrule
\end{tabular}
\end{table}


\subsection{Busca lexical: BM25}
O baseline lexical é o \textbf{BM25} \cite{robertson2009bm25}, com
\[
\mathrm{score}(q,d)=\sum_{t\in q}\mathrm{qtf}_t\cdot \ln\!\Big(1+\tfrac{N-\mathrm{df}_t+0{,}5}{\mathrm{df}_t+0{,}5}\Big)\cdot\frac{\mathrm{tf}_{t,d}\,(k_1+1)}{\mathrm{tf}_{t,d}+k_1\,(1-b+b\,|d|/\overline{|d|})}.
\]
Escolhemos o BM25 (e não o modelo vetorial TF-IDF/cosseno) porque a saturação do TF e a normalização por comprimento fazem diferença em resumos de tamanhos muito variados e porque é o baseline lexical padrão do BEIR, o que permite comparar os valores com a literatura. Os parâmetros $k_1=\BmK$ e $b=\BmB$ foram escolhidos no \emph{dev} (Tabela~\ref{tab:dev_bm25}); a superfície é plana e o ganho sobre os valores usuais é pequeno. O mesmo modelo lexical é usado em todos os experimentos.

\subsection{Tolerância a erros ortográficos}
\textbf{Correção.} Um termo da consulta é \emph{desconhecido} se for alfabético, tiver pelo menos três letras e não pertencer ao dicionário (tokens da coleção antes de stemming). Para cada termo desconhecido: (1) \emph{candidatos} são os termos do dicionário que compartilham pelo menos um k-grama (usando um índice k-grama$\rightarrow$termos, com marcadores de fronteira \texttt{\$}, $k=\SpK$); (2) o coeficiente de \emph{Jaccard} entre os conjuntos de k-gramas filtra os candidatos (limiar \SpJ); (3) a \emph{distância de edição} de Damerau-Levenshtein \cite{damerau1964,levenshtein1966} (inserção, remoção, substituição e troca de adjacentes custam 1; máximo de \SpEdit) seleciona o mais próximo, com desempate por Jaccard, frequência na coleção e ordem alfabética. Os parâmetros foram ajustados no \emph{dev} (Tabela~\ref{tab:dev_spell}).

\textbf{Protocolo do experimento.} Seguimos a padronização do enunciado: (i) dicionário antes do stemming; (ii) cada erro gera um termo \emph{inexistente} no dicionário (verificado por asserção); (iii) só entram consultas com NDCG@10 original $>0$ no sistema lexical; (iv) a recuperação por consulta é $\dfrac{\mathrm{NDCG}_{corr}-\mathrm{NDCG}_{erro}}{\mathrm{NDCG}_{orig}-\mathrm{NDCG}_{erro}}$ quando o denominador é positivo (senão registra-se que o erro não reduziu o desempenho). Para evitar escolha de casos favoráveis, as 10 consultas foram sorteadas (semente 42) entre as \TypoEligible{} consultas elegíveis do \emph{test}, sem usar qrels além do critério (iii). Cada consulta recebe \emph{um} erro, no termo de maior IDF (o mais informativo) com pelo menos 5 letras que não seja \emph{stopword}; a primeira letra é preservada; os quatro tipos de erro (substituição, inserção, remoção e troca de ordem) são atribuídos em rodízio. A correção é um experimento separado: os demais usam as consultas originais.

\subsection{Recuperação semântica}
Usamos um \emph{bi-encoder} \texttt{\EncName} (MiniLM, 6 camadas \cite{wang2020minilm}; arquitetura Sentence-BERT \cite{reimers2019sbert}) com saída de \EncDim{} dimensões, \emph{mean pooling} e normalização L2, de modo que o produto interno é o \textbf{cosseno}. Os embeddings dos documentos são calculados uma vez e armazenados; para cada consulta calcula-se o embedding e a similaridade com \emph{todos} os documentos (busca exaustiva, permitida pelo enunciado para esta coleção pequena). O comprimento máximo de entrada usado é \textbf{\EncMaxLen{} tokens}: os documentos têm em média \TokMean{} tokens (mediana \TokMedian{}; percentil 95: \TokPNinetyFive{}; máximo: \TokMax{}), logo \textbf{\EncTrunc{} documentos (\EncTruncPct\%) são truncados}. O limite foi escolhido no \emph{dev} junto com o modelo (Tabela~\ref{tab:dev_emb}).

\subsection{Busca híbrida: RRF}
Os 100 primeiros documentos dos rankings lexical e semântico são combinados por \emph{Reciprocal Rank Fusion} \cite{cormack2009rrf}: $\mathrm{RRF}(d)=\sum_{s}1/(k+\mathrm{rank}_s(d))$, com posição inicial 1; documentos ausentes de um ranking não recebem contribuição desse sistema, o que evita comparar scores de escalas diferentes (BM25 e cosseno). Empates são resolvidos pela melhor posição em algum ranking e, depois, pelo identificador. O parâmetro $k$ foi avaliado no \emph{dev} em \{5, 10, 20, 30, 40, 60, 80, 100, 150, 200\} (inclui o conjunto exigido \{10, 60, 100\}) e fixado em $k=\RrfK$ por maior NDCG@10 médio (Tabela~\ref{tab:dev_rrf}); o valor foi mantido no \emph{test}.

\subsection{Reranking}
O ranking híbrido é o recuperador inicial. Seus 50 primeiros documentos são reordenados pelo \emph{cross-encoder} \texttt{\CEName}, que codifica o par (consulta, documento) e produz um escore (logit) \cite{nogueira2019rerank}; os documentos abaixo da posição 50 permanecem na mesma ordem relativa do híbrido, de modo que existe um ranking completo para as métricas. O modelo foi escolhido no \emph{dev} (Tabela~\ref{tab:dev_ce}). O comprimento máximo do par é \CEMaxLen{} tokens; dos \CEPairs{} pares do \emph{test}, \CETruncPairs{} (\CETruncPct\%) foram truncados (média de \CEPairTokens{} tokens por par).

\subsection{RAG}
\textbf{Pipeline:} consulta $\rightarrow$ híbrido $\rightarrow$ reranking $\rightarrow$ seleção de contexto $\rightarrow$ LLM $\rightarrow$ resposta. Cada item do corpus é uma unidade de contexto (preserva a relação com os qrels); o contexto são os \textbf{Top-5} documentos pós-reranking, com título, texto e identificador (\texttt{[MED-\ldots]}). \textbf{Consultas:} as 10 consultas fixas do enunciado (semente 42 entre as 42 perguntas do \emph{test}); o código reproduz exatamente a lista do enunciado (verificado por asserção). \textbf{Condições:} (A) a LLM recebe somente a pergunta; (B) pergunta + Top-5 com seus identificadores, com instruções para responder \emph{apenas} com base no contexto, citar as fontes entre colchetes e declarar quando não houver evidência suficiente. Mesma LLM, mesmo idioma (inglês, o da coleção) e mesmo limite de 200 palavras nas duas condições. \textbf{LLM:} \LLMDescricao

\textbf{Avaliação.} Na recuperação: Recall@5 (denominador = total de relevantes), Precision@5 e Hit@5. Na geração: avaliação humana estruturada (0/1/2) de fidelidade, completude, utilidade e rastreabilidade, com planilha cega e embaralhada; na condição A, a fidelidade é julgada contra os mesmos Top-5 da condição B e a rastreabilidade é ``N/A'' (excluída da média).

\section{Avaliação}
\textbf{Dataset, \emph{splits} e qrels.} NFCorpus/BEIR: \NDocs{} documentos (campos \texttt{title} e \texttt{text}), \NQueries{} consultas e qrels de \emph{train}, \emph{dev} e \emph{test}. As consultas de cada \emph{split} são as que têm julgamentos no respectivo qrels (\NDev{} no \emph{dev}; \NTest{} no \emph{test}, com 12.334 julgamentos, graus 1 e 2). O \emph{dev} foi usado para escolher pré-processamento, $k_1$ e $b$, parâmetros da correção ortográfica, bi-encoder (modelo e limite de tokens), $k$ do RRF e cross-encoder; o \emph{test} só entrou na avaliação final, e seus qrels nunca modificaram consultas, parâmetros ou rankings. \textbf{Proveniência dos arquivos:} no ambiente de execução o servidor oficial do BEIR e o Hugging Face não estavam acessíveis; os arquivos foram obtidos de espelhos públicos no GitHub e \emph{validados}: duas cópias independentes são idênticas em documentos, consultas e qrels; as contagens batem com as do benchmark (3.633 documentos, 3.237 consultas, 323/324 consultas no \emph{test}/\emph{dev}, 42 perguntas no \emph{test}); e o código do enunciado reproduz exatamente as 10 consultas do RAG. O script \texttt{00\_download\_data.py} baixa o zip oficial e repete essas verificações.

\textbf{Métricas.} Precision@10, Recall@10, MAP@10 e NDCG@10, calculadas por consulta e depois promediadas. Para P, R e AP, relevante é todo documento com grau $>0$; Recall e AP dividem pelo \emph{total} de relevantes da consulta nos qrels; AP@10 soma $P@k$ nas posições $k\le10$ com relevante; NDCG@10 usa ganho linear (o grau) e ranking ideal calculado com todos os julgamentos da consulta. As quatro métricas foram implementadas no projeto e conferidas contra a \texttt{ir\_measures} \cite{macavaney2022irmeasures} nos quatro sistemas do \emph{test}: a maior diferença absoluta por consulta foi $\le 10^{-15}$. Para comparar sistemas usamos o teste de Wilcoxon pareado sobre o NDCG@10 por consulta e intervalos de confiança \emph{bootstrap}.

\textbf{Modelos e parâmetros finais.} Lexical: BM25 ($k_1=\BmK$, $b=\BmB$; stopwords removidas; stemmer \PreStem). Semântico: \texttt{\EncName}, máximo de \EncMaxLen{} tokens, cosseno. Híbrido: RRF com $k=\RrfK$ sobre os 100 primeiros de cada sistema. Reranking: \texttt{\CEName} sobre o Top-50 híbrido.

\textbf{Ambiente computacional.} Todos os experimentos de recuperação rodaram em CPU: \CpuModel, \CpuCores{} núcleo(s), \RamGB{} GB de RAM, sem GPU; Python \PyVer, NumPy \NumpyVer, ONNX Runtime \OrtVer, NLTK \NltkVer. Como o ambiente não tinha PyTorch nem acesso ao Hugging Face, os \emph{mesmos pesos} dos modelos foram executados com ONNX Runtime (bi-encoder) e com uma implementação própria do encoder BERT em NumPy (cross-encoder), ambas validadas: os embeddings do bi-encoder reproduzem os gerados pelo \texttt{sentence-transformers} (cosseno $=1{,}000000$ em 64 documentos) e o cross-encoder reproduz os escores de referência publicados no card do modelo ($8{,}607141$ e $-4{,}320078$ contra $8{,}607138$ e $-4{,}320078$). O código também traz os \emph{backends} do \texttt{sentence-transformers} (padrão no Colab/PC), que não foram executados neste ambiente. Os tempos reportados valem para este hardware.

\textbf{Protocolo.} (1) ajustes no \emph{dev}; (2) congelamento das escolhas (\texttt{results/dev\_selection.json}); (3) execução única no \emph{test} com os mesmos Top-100 por sistema salvos em CSV; (4) análises por consulta com seleção automática dos casos (por diferença de NDCG@10), sem escolha manual.
